\chapter*{Introduction}
\addstarredchapter{Introduction}
%%%%%%%%%%%%%%%%%%%%%%%%%%%%%%%%%%%%%%%%%%%%%%%%%%%%%%%%%%%%%%%%%%%%%%%%%%%%%%%%%%%%%%%%%%%%%%%%%%%%%%%%%%%%%%%%%%%%%%%%%%%%%%%%
 Au cours des dernières années, le développement rapide des technologies numériques a profondément transformé de nombreux secteurs, favorisant la digitalisation des processus et l’amélioration de l’expérience utilisateur. Le domaine de la location de véhicules s’inscrit pleinement dans cette dynamique, où l’intégration de solutions numériques innovantes constitue un levier essentiel pour optimiser la gestion des agences, améliorer la satisfaction des clients et assurer un meilleur suivi des activités.
 
 Cependant, malgré ces avancées, la gestion des opérations de location au sein de plusieurs agences demeure souvent complexe lorsqu’elle repose sur des méthodes manuelles ou des systèmes non centralisés. Cette situation peut engendrer des erreurs, des retards ainsi que des pertes d’informations, notamment dans le suivi des clients, des réservations et des transactions financières. Par conséquent, les agences de location éprouvent un besoin croissant de disposer d’une solution centralisée, fiable et performante permettant de gérer efficacement l’ensemble du processus de location.
 
 Dans ce contexte, et dans le cadre de l’obtention de notre diplôme de licence en Business Computing, spécialité E-Business, à l’\textbf{École Supérieure d’Économie Numérique}, nous avons été amenés, lors de notre stage au sein de la société \textbf{Yes Internet}, à concevoir et développer une \textbf{plateforme web centralisée de gestion de location de véhicules}. Cette solution vise à digitaliser les processus métier, faciliter la gestion des véhicules, des clients et des réservations, assurer un suivi rigoureux des paiements, ainsi qu’à offrir des outils décisionnels à travers des tableaux de bord interactifs.

Ce rapport est structuré en plusieurs chapitres :
\begin{itemize}
    \item \textbf{Chapitre 1 – Présentation du projet} : expose le contexte général, analyse les limites des méthodes actuelles et justifie la nécessité d’une solution numérique adaptée aux besoins des agences.
    \item \textbf{Chapitre 2 – Planification et architecture} : décrit l’identification des besoins fonctionnels, la modélisation des cas d’utilisation et l’organisation du projet en sprints de développement.
    \item \textbf{Chapitre 3 – Sprint 1 : Gestion des intervenants} : présente le déroulement du premier sprint, incluant la planification et le développement des fonctionnalités liées à la gestion des intervenants.
    \item \textbf{Chapitre 4 – Sprint 2 : Gestion des processus des réservations} : détaille la mise en œuvre des fonctionnalités liées aux procéssus réservations.
    \item \textbf{Chapitre 5 – Phase de clôture} : met en lumière les outils et technologies et l'architecture utilisés pour le développement de la plateforme.
\end{itemize}

Le rapport se termine par des perspectives d'amélioration, notamment l'intégration d'un gateway de paiement en ligne, le suivi du comportement des clients et le développement d'une application mobile.



